%% ============================================================
%% CHAPITRE 4 — CONCEPTION ET DEVELOPPEMENT DE LA PARTIE LOGICIELLE
%% ============================================================
%% Session  : 4
%% Auteur   : [Nom etudiant]
%% Redige le: 12 mai 2026
%% Inspecte : orch/+orch/run.m  (sequencement des phases)
%%            orchestrator/cli.py (CLI Python : argparse, MATLAB -batch)
%%            doc_generator/{cli.py, consolidate.py, prompt.py, llm_client.py}
%%            analyzer/python/analyzer/{engine.py, rules.py, facts.py,
%%              reporter.py, cli.py}
%%            analyzer/+ana/collect.m + analyzer/knowledge/*.yaml
%%            validator/{+val/, mil/, sil/, pil/, cov/, plant/, +fp/}
%%            app/bms_assistant.m (UI 5 vues : Dashboard / Validate /
%%              Reports / Plots / Model)
%%            tests/python/ (8 fichiers, 42 tests)
%%            tests/matlab/ (3 fichiers)
%% ============================================================

\chapter{Conception et développement de la partie logicielle}
\label{chap:partie_logicielle}

\section{Introduction}
\label{sec:ch4_introduction}

Le chapitre~\ref{chap:partie_modele} a clos la branche Modèle du
V-model dual en livrant un système sous test fonctionnel et validé
unitairement~: \emph{plant} ECM~2RC, BMS \emph{master} et \emph{slave},
prédicteur LSTM intégré au \emph{master}. Le présent chapitre déroule
symétriquement la branche Logicielle. L'objectif n'est plus de
modéliser la batterie ou le calculateur, mais de construire l'outil
intelligent qui pilote, supervise et synthétise la chaîne de validation
décrite au chapitre~3.

L'enjeu est double. D'une part, automatiser une chaîne hétérogène qui
mélange du code MATLAB de validation, du code Python d'analyse et un
service externe d'inférence neuronale (\textbf{API Groq}) en un
pipeline unique reproductible et instrumenté. D'autre part, exposer
cette chaîne à un ingénieur de validation BMS sous une forme
ergonomique et sans courbe d'apprentissage, par une interface
graphique \textsc{MATLAB App~Designer} intégrée à l'environnement
quotidien de l'équipe.

Le chapitre suit la même structure descendante~\textit{$\to$}~ascendante
que le précédent. Les exigences logicielles héritées du
chapitre~\ref{chap:analyse_specification} et l'architecture modulaire
sont d'abord posées (section~\ref{sec:ch4_exigences_archi}). Les
quatre modules logiciels — Validateur, Analyseur, Générateur de
rapports et Orchestrateur — sont ensuite développés en
section~\ref{sec:ch4_modules}, suivis de l'interface graphique
(section~\ref{sec:ch4_ui}). Le chapitre se conclut par la validation
unitaire du logiciel (section~\ref{sec:ch4_validation_unitaire}) et
par une transition vers la phase d'intégration et d'acceptation,
traitée au chapitre~\ref{chap:integration_validation}.

\section{Exigences et architecture logicielle}
\label{sec:ch4_exigences_archi}

\subsection{Exigences logicielles héritées}
\label{subsec:ch4_exigences}

La branche Logicielle s'attache à satisfaire les vingt-neuf
exigences \texttt{REQ-SW-*} cataloguées au
chapitre~\ref{chap:analyse_specification}. Elles se répartissent en
cinq familles correspondant aux cinq composants logiciels du
système. La table~\ref{tab:ch4_familles_req_sw} résume cette
répartition et précise pour chaque famille le module responsable
de la satisfaction.

\begin{table}[H]
  \centering
  \caption{Répartition des 29~exigences logicielles
  \texttt{REQ-SW-*} par module responsable.}
  \label{tab:ch4_familles_req_sw}
  \small
  \begin{tabularx}{\textwidth}{|l|c|X|}
    \hline
    \textbf{Préfixe} & \textbf{N} & \textbf{Module responsable} \\
    \hline
    \texttt{REQ-SW-VAL-*}  & 7 & Validateur (orchestration MIL/SIL/PIL/COV) \\
    \hline
    \texttt{REQ-SW-ANA-*}  & 6 & Analyseur (collecte facts + moteur de règles) \\
    \hline
    \texttt{REQ-SW-DOC-*}  & 6 & Générateur de rapports (consolidation + LLM) \\
    \hline
    \texttt{REQ-SW-ORCH-*} & 5 & Orchestrateur (séquencement, CLI, callbacks) \\
    \hline
    \texttt{REQ-SW-UI-*}   & 5 & Interface utilisateur (App~Designer) \\
    \hline
    \textbf{Total}        & \textbf{29} & — \\
    \hline
  \end{tabularx}
\end{table}

\subsection{Architecture modulaire générale}
\label{subsec:ch4_archi_modulaire}

L'architecture logicielle adopte le \emph{principe de séparation des
responsabilités} en isolant chaque préoccupation dans un module
indépendant, déployable et testable séparément. La
figure~\ref{fig:ch4_archi_modulaire} présente cette architecture en
cinq blocs~: un \textbf{Validateur} qui exerce le système sous test
en MIL, SIL et PIL et calcule la couverture structurelle, un
\textbf{Analyseur} qui consomme les artefacts du Validateur et y
applique un moteur expert pour produire un diagnostic, un
\textbf{Générateur de rapports} qui consolide ces sorties et les
transforme en rapport HTML lisible par l'API Groq, un
\textbf{Orchestrateur} qui pilote les trois précédents en respectant
leurs dépendances, et une \textbf{Interface utilisateur} qui expose
l'ensemble à l'ingénieur de validation BMS.

\begin{figure}[H]
  \centering
  \includegraphics[width=0.95\textwidth]{diagrams/ch4_archi_modulaire.png}
  \caption{Architecture modulaire générale de la partie logicielle.
  Chaque module est indépendant et testable séparément~; les flèches
  représentent les flux de données et de contrôle.}
  \label{fig:ch4_archi_modulaire}
\end{figure}

Cette organisation présente trois propriétés clés vis-à-vis des
exigences. La \emph{réversibilité} du choix de service LLM est
garantie par l'isolation de l'appel HTTP dans un client minimal
(\texttt{doc\_generator/llm\_client.py}, 100~lignes) qui parle le
protocole \emph{OpenAI-compatible}~: trois variables d'environnement
(\texttt{LLM\_API\_URL}, \texttt{LLM\_MODEL}, \texttt{LLM\_API\_KEY})
suffisent pour basculer entre Groq, WisGate, OpenAI ou
OpenRouter sans recompiler. La \emph{traçabilité de bout en bout} est
assurée par la transmission, à chaque étage, d'un identifiant de run
horodaté (\texttt{ts}) qui suffixe tous les artefacts produits. La
\emph{reproductibilité} est obtenue par l'écriture systématique d'un
\emph{snapshot} des paramètres et des sorties dans
\texttt{validator/reports/} et \texttt{analyzer/reports/}.

\subsection{Architecture détaillée~: flux de données inter-modules}
\label{subsec:ch4_archi_detaillee}

Les modules communiquent exclusivement par fichiers JSON normalisés,
ce qui découple totalement les langages d'implémentation et autorise
l'inspection manuelle de tout étage intermédiaire. La
figure~\ref{fig:ch4_flux_inter_modules} décrit les flux entre les
fichiers d'artefacts produits ou consommés à chaque étage, avec leurs
schémas et leurs chemins relatifs au dépôt.

\begin{figure}[H]
  \centering
  \includegraphics[width=0.95\textwidth]{diagrams/ch4_flux_inter_modules.png}
  \caption{Diagramme de flux inter-modules~: artefacts JSON et HTML
  échangés entre Validateur, Analyseur, Générateur de rapports,
  Orchestrateur et Interface utilisateur.}
  \label{fig:ch4_flux_inter_modules}
\end{figure}

\subsection{Diagramme de classes principal}
\label{subsec:ch4_diag_classes}

Le diagramme UML de classes (figure~\ref{fig:ch4_diag_classes}) ne
représente que les composants logiciels développés dans le cadre du
projet, à l'exclusion des bibliothèques tierces (\textsc{torch},
\textsc{numpy}, \texttt{urllib.request}, etc.). La hiérarchie est
volontairement plate~: aucun héritage profond, aucune
métaprogrammation, conformément à la règle interne du projet selon
laquelle le code Python doit rester lisible par un ingénieur
familier de \textsc{MATLAB} sans formation à la programmation
orientée objet avancée.

\begin{figure}[H]
  \centering
  \includegraphics[width=0.95\textwidth]{diagrams/ch4_diag_classes.png}
  \caption{Diagramme UML de classes principal de la partie logicielle.
  La classe \texttt{bms\_assistant} hérite de
  \texttt{matlab.apps.AppBase}~; les autres composants sont des modules
  fonctionnels Python ou des packages \textsc{MATLAB} (\texttt{+orch},
  \texttt{+ana}, \texttt{+val}).}
  \label{fig:ch4_diag_classes}
\end{figure}

\subsection{Diagramme de séquence~: scénario de validation complet}
\label{subsec:ch4_diag_seq}

La figure~\ref{fig:ch4_diag_seq} décrit le scénario complet d'une
exécution de validation tel que déclenché par l'ingénieur depuis la
vue \emph{Validate} de l'interface graphique. Ce scénario, qui
correspond au cas d'utilisation \texttt{UC-VAL-01} défini au
chapitre~\ref{chap:analyse_specification}, exerce successivement les
huit phases de la chaîne~: chargement des paramètres (\texttt{init\_system.m}),
MIL, SIL, PIL, COV, collecte des \emph{facts} (\texttt{ana.collect}),
analyseur Python, et génération des deux rapports HTML.

\begin{figure}[H]
  \centering
  \includegraphics[width=0.95\textwidth]{diagrams/ch4_diag_sequence.png}
  \caption{Diagramme UML de séquence du scénario de validation
  complet (\texttt{UC-VAL-01}). Le \emph{progress\_cb} remonte
  l'avancement à l'UI à chaque test individuel.}
  \label{fig:ch4_diag_seq}
\end{figure}

\section{Développement des modules logiciels}
\label{sec:ch4_modules}

\subsection{Module Validateur}
\label{subsec:ch4_validateur}

Le \textbf{Validateur} est le module qui exerce effectivement le
système sous test du chapitre~\ref{chap:partie_modele}. Il est
implémenté en \textsc{MATLAB} sous la forme de cinq packages~:
\texttt{+plant} pour la convergence du \emph{plant model},
\texttt{+mil} pour les tests \emph{Model-in-the-Loop},
\texttt{+sil} pour les tests \emph{Software-in-the-Loop},
\texttt{+pil} pour les tests \emph{Processor-in-the-Loop} sur la cible
STM32H7A3, et \texttt{+cov} pour la collecte de la couverture
structurelle. Un sixième package \texttt{+val} regroupe les
\emph{helpers} partagés (résolution des signaux loggés, génération du
\texttt{Simulink.SimulationData.Dataset} d'entrée du \emph{master},
sérialisation JSON des verdicts).

\paragraph{Validation de la convergence du \emph{plant model}.}
La fonction \texttt{plant.run()} compare, sur les soixante-dix
trajets disponibles, les sorties simulées
(\(V_\text{pack}\), \(\mathrm{SoC}_\text{pack}\), \(T_\text{pack}\))
aux mesures correspondantes. Les RMSE par trajet sont sérialisées
dans \path{validator/reports/plant_convergence.json}, et les
verdicts globaux (médiane, moyenne, max) sont calculés en regard des
seuils \texttt{req.PLT\_*\_max\_*} chargés depuis
\texttt{init\_system.m}. Cette fonction satisfait
\texttt{REQ-SW-VAL-PLT-01}.

\paragraph{Validation automatisée MIL / SIL / PIL.}
Les trois fonctions \texttt{mil.run()}, \texttt{sil.run()} et
\texttt{pil.run()} suivent le même schéma~: découverte des fichiers
\texttt{test\_<REQ-ID>.m} dans le sous-dossier \texttt{tests/} de
chaque package, exécution sous \emph{harness} \textsc{Simulink},
calcul des verdicts, écriture du rapport JSON horodaté avec une copie
\emph{sticky} non horodatée pour faciliter les comparaisons. La
fonction \texttt{sil.run()} peut sauter la régénération du code
(\texttt{skip\_build}) lorsque le code C n'a pas changé depuis le
dernier \emph{run}, ce qui économise environ 30~secondes par
exécution. La fonction \texttt{pil.run()} déclenche, si nécessaire, la
recompilation du binaire pour la carte
\emph{NUCLEO-H7A3ZI-Q} via la chaîne \emph{Embedded Coder} configurée
au chapitre~\ref{chap:partie_modele}, opération qui prend
typiquement de l'ordre d'une minute par phase. Ces trois fonctions
satisfont collectivement
\texttt{REQ-SW-VAL-MIL-01}, \texttt{REQ-SW-VAL-SIL-01} et
\texttt{REQ-SW-VAL-PIL-01}.

\paragraph{Analyse de la couverture et détection du code mort.}
La couverture structurelle est calculée \emph{en ligne} avec les
tests MIL, et non en passe séparée~: le \emph{driver} \texttt{mil.run}
collecte un objet \texttt{cvdata} pendant son exécution, que la
fonction \texttt{cov.from\_cvdata()} agrège ensuite par modèle
(\emph{slave}, \emph{master}, \emph{predictor}) et par métrique
(\emph{Execution}, \emph{Decision}, \emph{Condition}, \emph{MC/DC}).
Cette stratégie en \emph{passe unique} évite de re-simuler la totalité
de la suite MIL, ce qui réduit le temps total d'environ 40~\%, et
garantit par construction que la couverture mesurée correspond
exactement au stimulus exercé. Le verdict est binaire~: une exigence
\texttt{REQ-LL-COV-DEAD-01} est satisfaite si et seulement si zéro
bloc reste inatteignable au sens de la métrique \emph{Execution}, qui
constitue le \emph{proxy} canonique de code mort au niveau modèle.
Les autres métriques sont remontées à titre informationnel.

\begin{figure}[H]
  \centering
  \fbox{\parbox{0.85\textwidth}{\centering\vspace{2.2cm}
    \textit{[Capture de la sortie console de cov.run : tableau
    récapitulatif par modèle (slave, master, predictor) avec les 4
    métriques (exec/dec/cond/mcdc) et leur pourcentage de couverture
    + nombre de blocs unhit, + verdict PASS]}
    \vspace{2.2cm}}}
  \caption{Sortie de couverture structurelle produite par
  \texttt{cov.from\_cvdata()} pour la dernière campagne MIL.}
  \label{fig:ch4_cov_output}
\end{figure}

\subsection{Module Analyseur}
\label{subsec:ch4_analyseur}

Le \textbf{Analyseur} est implémenté en deux parties complémentaires~:
un \emph{collecteur de faits} \textsc{MATLAB}
(\texttt{analyzer/+ana/}) qui extrait les métriques statiques des
modèles \textsc{Simulink}, et un \emph{moteur expert} \textsc{Python}
(\texttt{analyzer/python/analyzer/}) qui applique un ensemble de
règles déclaratives à ces métriques pour produire un diagnostic
catégorisé.

\paragraph{Collecte de métriques par analyse statique.}
La fonction \texttt{ana.collect} parcourt les modèles
\texttt{bms\_master}, \texttt{bms\_slave} et
\texttt{fault\_predictor} et y applique huit sondes statiques
indépendantes~: \texttt{metrics\_static} (compte de blocs, profondeur
hiérarchique, complexité cyclomatique), \texttt{metrics\_stateflow}
(nombre d'états, transitions, jonctions par chart),
\texttt{metrics\_solver} (paramètres de simulation), \texttt{metrics\_codegen}
(état de la dernière génération de code C), \texttt{metrics\_signals}
(signaux nommés, types, plages),
\texttt{metrics\_rates} (taux d'échantillonnage et hiérarchie),
\texttt{metrics\_data} (variables de modèle et bus), et
\texttt{metrics\_codeblock} (sondes spécifiques aux blocs
\emph{MATLAB Function} et \emph{Simulink Function}). Le résultat est
sérialisé en \texttt{analyzer/reports/FACTS\_<timestamp>.json} avec
un schéma versionné (\texttt{schema\_version: "1.0"}) qui isole le
moteur expert des évolutions futures du collecteur. Cette fonction
satisfait \texttt{REQ-SW-ANA-FACTS-01}.

\paragraph{Système expert~: règles d'inférence en YAML.}
Le moteur expert applique une logique de
\emph{forward chaining}~\cite{Forsberg1991} sur une base de règles
déclaratives au format YAML, stockées sous
\texttt{analyzer/knowledge/}. Chaque règle décrit une condition
(\texttt{when}), un ensemble de \emph{bindings} (\texttt{bind}), une
sévérité (\texttt{severity} \(\in\) \{info, minor, major, critical\}),
un message à émettre (\texttt{finding}), un raisonnement
(\texttt{reasoning}), une liste de recommandations
(\texttt{recommend}) et des références normatives (\texttt{refs}). Le
moteur (\texttt{engine.py}) charge ces règles, les pré-compile en AST
Python par \texttt{ast.parse}, puis les évalue dans un interprète
restreint qui n'autorise que les opérations de comparaison,
combinaison logique et résolution de chemins pointés sur la base de
faits. Toute valeur manquante évalue à \emph{faux}~: la règle ne
déclenche simplement pas, sans erreur. Cette discipline garantit
qu'un fichier de \emph{facts} incomplet (par exemple lorsqu'un module
n'a pas été collecté) ne fait pas échouer l'analyse mais réduit
silencieusement son périmètre. La conception satisfait
\texttt{REQ-SW-ANA-RULES-01}.

\paragraph{Politique ASIL~: \texttt{asil\_policy.yaml}.}
Une règle de modélisation interne au projet stipule que les seuils
chiffrés (par exemple \emph{\og{}le nombre de blocs morts ne doit pas
dépasser X~\%\fg}) ne doivent jamais être codés en dur dans les
règles, mais centralisés dans un unique fichier
\texttt{asil\_policy.yaml}. Ce fichier traduit le contexte ASIL D du
projet en seuils numériques par modèle et par métrique, et constitue
le seul point de réglage si l'on souhaite resserrer ou relâcher les
exigences pour une variante du système. La conformité au standard est
garantie par les références \texttt{refs} attachées à chaque règle,
qui citent explicitement la clause d'ISO~26262-6:2018
correspondante~\cite{ISO26262_6_2018} (par exemple Table~12 pour les
métriques de couverture structurelle). Conformément à une décision
prise dès l'analyse, ISO~26262-8 (qualification des outils) n'est
\emph{jamais} citée pour des sujets de style de modélisation~: cette
norme est réservée à la qualification éventuelle de l'outil lui-même.

\paragraph{\emph{Downrank} des sous-systèmes en cours de développement.}
Le moteur dispose d'une logique dite \emph{WIP downrank}~: lorsqu'un
\emph{finding} concerne un bloc dont le chemin tombe dans la liste
\texttt{policy.wip\_paths}, sa sévérité est rétrogradée à \emph{info}
et le \emph{finding} est marqué d'un drapeau \texttt{wip:~true}. Cela
permet de ne pas voir le tableau de bord saturé d'avertissements sur
des sous-systèmes encore en cours de développement, tout en
préservant la traçabilité complète de l'analyse. Le \emph{locator}
\textsc{MATLAB} (\texttt{hilite\_system('chemin/du/bloc')})
généré pour chaque \emph{finding} non trivial permet à l'ingénieur
de sauter directement au bloc fautif depuis le rapport.

\paragraph{Génération du rapport d'analyse.}
Le \texttt{reporter.py} sérialise le résultat en
\path{analyzer/reports/ANALYSIS_<timestamp>.json} avec, en
parallèle, une copie sticky \texttt{ANALYSIS.json} pour faciliter les
comparaisons entre exécutions consécutives. Un \emph{digest} textuel
est imprimé sur la sortie standard, regroupé par sévérité décroissante
(\emph{critical} d'abord), avec le \emph{locator} et la première
recommandation pour chaque \emph{finding}. Cette présentation
condensée permet à l'ingénieur d'agir immédiatement sans ouvrir le
fichier JSON complet.

\begin{figure}[H]
  \centering
  \fbox{\parbox{0.85\textwidth}{\centering\vspace{2.2cm}
    \textit{[Exemple de digest console de l'analyseur : section
    CRITICAL avec 1 finding (DC-MASTER-EXEC : execution dead = X.YY\%
    > limit Z\%), section MAJOR vide, section MINOR avec 2-3 findings,
    INFO avec WIP tag visible. Verdict global PASS/FAIL final.]}
    \vspace{2.2cm}}}
  \caption{Exemple de \emph{digest} produit par l'analyseur~: regroupement
  par sévérité décroissante, \emph{locator} \textsc{MATLAB} et première
  recommandation par \emph{finding}.}
  \label{fig:ch4_analyzer_digest}
\end{figure}

\subsection{Module Générateur de rapports}
\label{subsec:ch4_doc_generator}

Le \textbf{Générateur de rapports} (\texttt{doc\_generator/}) est le
module qui transforme les artefacts JSON techniques produits par les
deux modules précédents en un document de synthèse lisible par un
ingénieur ou un manager. Il s'organise en quatre étapes qui
correspondent à autant de fichiers~: \texttt{consolidate.py}
(consolidation), \texttt{prompt.py} (composition du prompt LLM),
\texttt{llm\_client.py} (client HTTP minimal) et \texttt{cli.py}
(orchestration de l'ensemble).

\paragraph{Consolidation des résultats.}
Le \texttt{consolidate.py} construit deux artefacts à partir des
quatre rapports \texttt{MIL.json}, \texttt{SIL.json}, \texttt{PIL.json}
et \texttt{ANALYSIS.json}~: un \texttt{FULL\_REPORT.json} qui
concatène les quatre sources \emph{sans aucune modification} sous des
clés stables, et un \texttt{slim} dérivé qui sert d'entrée au LLM. Le
\texttt{FULL\_REPORT.json} est conservé tel quel parce qu'il est
référencé par \emph{JSON Pointer} (RFC~6901) depuis le rapport HTML
final~: chaque assertion du rapport HTML pointe ainsi vers la ligne
exacte du fichier brut qui la justifie, ce qui rend le rapport
auditable bloc par bloc. Cette propriété d'auditabilité satisfait
\texttt{REQ-SW-DOC-AUDIT-01}.

\paragraph{Génération du rapport résumé en deux parties.}
Le rapport HTML est produit en \emph{deux} fichiers distincts plutôt
qu'en un seul, pour une raison purement opérationnelle~: le quota
\emph{Token-Per-Minute} (TPM) du palier gratuit de l'\textbf{API Groq}
plafonne à 8000~\emph{tokens} par appel sur le modèle
\texttt{openai/gpt-oss-120b} (\emph{prompt}~+~\emph{max\_completion}
confondus). En scindant le rapport en deux appels — Partie~1
(\og{}Validation Data\fg) et Partie~2 (\og{}Engineering Analysis\fg) —
chaque appel dispose de son propre budget, ce qui multiplie par deux
la marge utilisable. Le \texttt{slim} consommé par chaque appel est
construit par les fonctions \texttt{slim\_for\_part1} et
\texttt{slim\_for\_part2} qui ne retiennent que les champs strictement
nécessaires à la section ciblée, et qui marquent chaque ligne avec son
JSON Pointer dans le \texttt{FULL\_REPORT.json}. Une étape finale
de \emph{trim} progressif (\texttt{\_trim\_part1},
\texttt{\_trim\_part2}) garantit que la somme \emph{prompt}~+~\emph{max\_completion}
reste sous le budget cible \texttt{TPM\_BUDGET} = 7600~\emph{tokens},
qui laisse environ 5~\% de marge sur le plafond TPM réel.

\paragraph{Client LLM minimal et portable.}
Le \texttt{llm\_client.py} (\texttt{$<$~100 lignes}) parle le
protocole \emph{OpenAI-compatible} via la \emph{stdlib} Python
exclusivement (\texttt{urllib.request} + \texttt{json}), ce qui évite
toute dépendance à \texttt{openai-python} ou à \texttt{requests} et
réduit la surface d'attaque comme l'empreinte d'installation. Le
choix du fournisseur (Groq, WisGate, OpenAI, OpenRouter) est piloté
par les variables d'environnement \texttt{LLM\_API\_URL},
\texttt{LLM\_MODEL} et \texttt{LLM\_API\_KEY}~; aucune logique
fournisseur-spécifique n'apparaît dans le code. Trois précautions
opérationnelles sont implémentées explicitement~: l'envoi d'un
\texttt{User-Agent} explicite (sans quoi le proxy Cloudflare devant
\texttt{api.groq.com} renvoie un HTTP~403), l'envoi simultané des
deux clés \texttt{max\_tokens} et \texttt{max\_completion\_tokens}
(pour rester compatible avec OpenAI \emph{strict}), et le passage
\texttt{reasoning\_effort: "low"} qui plafonne les \emph{tokens} de
raisonnement cachés du modèle \texttt{gpt-oss-120b} à environ 100,
contre 1200 par défaut. Une logique de \emph{retry} unique sur erreur
HTTP transitoire (429, 500, 502, 503, 504) couvre l'attente d'environ
60~secondes nécessaire pour que la fenêtre glissante TPM se libère.

\paragraph{Audit des références.}
Une étape de post-traitement (\texttt{\_audit\_refs}) extrait du HTML
généré tous les \texttt{<code>/...</code>}, les résout dans le
\texttt{FULL\_REPORT.json} via le résolveur RFC~6901 maison, et liste
les pointeurs cassés. Cette étape protège contre les hallucinations
du modèle qui inventerait un chemin JSON inexistant~: tout pointeur
non résolu est signalé et bloque, le cas échéant, la publication du
rapport.

\begin{figure}[H]
  \centering
  \fbox{\parbox{0.85\textwidth}{\centering\vspace{2.2cm}
    \textit{[Capture d'un extrait du rapport HTML généré : section
    SUMMARY avec table KPI (tests total, passed, failed, MIL/SIL/PIL
    pass-rate, suites run), section MIL avec table req\_id|name|
    status|reference où chaque ref est un <code>/...</code> cliquable
    pointant vers FULL\_REPORT.json. Style sans-serif noir/blanc.]}
    \vspace{2.2cm}}}
  \caption{Extrait du rapport HTML produit par le Générateur~: chaque
  référence \texttt{<code>} est un pointeur RFC~6901 vers le
  \texttt{FULL\_REPORT.json} consolidé.}
  \label{fig:ch4_doc_html}
\end{figure}

\subsection{Module Orchestrateur}
\label{subsec:ch4_orchestrateur}

L'\textbf{Orchestrateur} est le module qui coordonne les trois
modules précédents et qui sert d'intermédiaire entre l'interface
utilisateur \textsc{MATLAB} et la chaîne d'exécution. Il est
implémenté \emph{deux fois}~: une version \textsc{MATLAB}
(\texttt{orch/+orch/run.m}) consommée par l'interface graphique et
optimisée pour la transmission temps réel des callbacks de
progression, et une version \textsc{Python} (\texttt{orchestrator/cli.py})
consommée en mode \emph{batch} pour l'intégration dans les pipelines
CI/CD. Cette double implémentation répond à des contraintes
d'environnement~: l'UI ne peut pas appeler \texttt{matlab -batch} en
sous-processus depuis elle-même, et le pipeline CI ne dispose pas
toujours d'une session \textsc{MATLAB} interactive.

\paragraph{Coordination et séquencement des phases.}
La fonction \texttt{orch.run} construit un plan d'exécution
(\texttt{local\_plan}) sous forme de liste ordonnée de phases, dont
chacune est un \texttt{struct} contenant un nom (\texttt{name}), un
pointeur de fonction (\texttt{fcn}), un poids relatif
(\texttt{weight}) servant à pondérer la barre de progression, et un
nombre de tests prévus (\texttt{test\_count}). Les poids ne sont pas
arbitraires~: ils reflètent la durée empirique relative observée sur
plusieurs campagnes (un test PIL coûte typiquement 15~unités, un
test MIL 4~unités, le \emph{build} SIL 30~unités, le \emph{build}
PIL 60~unités). Cette pondération permet à la barre de progression
de l'UI de progresser linéairement en temps réel, sans à-coups, ce
qui est crucial pour la perception qu'a l'utilisateur de la
\emph{réactivité} de l'outil.

\paragraph{Gestion des flux et des données d'entrée.}
Chaque phase reçoit en entrée le \texttt{struct opt} consolidé par
l'\texttt{inputParser} (filtres MIL/SIL/PIL, \emph{flags}
\texttt{skip\_*}, callbacks de \emph{logging} et de progression),
le chemin racine du dépôt et deux callbacks~: \texttt{log\_cb} pour
la sortie console et \texttt{progress\_cb} pour la mise à jour de
la barre. Les phases produisent un \texttt{struct rep} qui est
agrégé dans le résultat final \texttt{out.paths.<phase>}. Cette
uniformité d'interface entre les six phases permet à
\texttt{orch.run} de les enchaîner sans connaître leur sémantique
interne.

\paragraph{Gestion des états et des erreurs.}
La machine à états de l'orchestrateur est volontairement simple~:
chaque phase échoue indépendamment et la première erreur arrête le
pipeline en remontant à l'UI le nom de la phase fautive
(\texttt{out.failed\_phase}) et le message
(\texttt{out.error}). Aucun \emph{rollback} automatique n'est tenté~:
les artefacts produits par les phases déjà terminées sont conservés
sur disque, ce qui permet à l'ingénieur de les inspecter ou de
relancer manuellement les phases ratées en passant
\texttt{skip\_<phase>} à la prochaine invocation. Cette philosophie
\emph{fail-fast~+~conserve} a été préférée à un mécanisme de
transactions plus complexe parce qu'elle correspond mieux aux
habitudes des ingénieurs de validation, qui préfèrent re-jouer une
phase isolée que reprendre la totalité du pipeline.

\paragraph{Communication avec l'interface utilisateur.}
La communication entre l'UI et l'orchestrateur s'effectue par
\emph{callbacks} synchrones plutôt que par \emph{events}
asynchrones~: la fonction \texttt{orch.run} appelle directement
\texttt{progress\_cb(name, fraction)} et
\texttt{log\_cb(message)} à chaque transition d'intérêt, et
l'UI met à jour ses \emph{phase chips} et son log dans le même
\emph{tick} de la boucle d'événements \textsc{MATLAB}. Ce choix simplifie
considérablement la gestion des erreurs (pas de file d'événements
à drainer en cas d'exception) au prix d'un couplage modéré entre
l'UI et l'orchestrateur, jugé acceptable compte tenu du fait que les
deux composants vivent dans le même processus.

\begin{figure}[H]
  \centering
  \includegraphics[width=0.85\textwidth]{diagrams/ch4_diag_etat_orch.png}
  \caption{Diagramme UML d'états de l'Orchestrateur. Chaque phase
  \emph{Running} émet à intervalle régulier un événement
  \texttt{progress\_cb} consommé par la barre de progression de l'UI.}
  \label{fig:ch4_diag_etat_orch}
\end{figure}

\section[Développement de l'interface graphique]{Développement de l'interface graphique (\textsc{MATLAB App~Designer})}
\label{sec:ch4_ui}

\subsection{Choix technologique et contraintes}
\label{subsec:ch4_choix_techno_ui}

L'interface graphique est développée en \textsc{MATLAB
App~Designer}, choix dicté par trois contraintes~: l'absence de
nouvelle dépendance pour l'utilisateur final (App~Designer est
fourni avec toute installation \textsc{MATLAB}), la possibilité
d'appeler directement les fonctions du package \texttt{+orch} sans
bridge inter-processus, et l'intégration native avec
\textsc{Simulink} pour la fonction d'inspection des modèles
\emph{plant} et SUT. Le revers de la médaille est l'expressivité
limitée de la grille \texttt{uigridlayout} comparée à un framework
\emph{web} moderne, qui impose une discipline de mise en page plus
stricte. Le choix a néanmoins été retenu parce que l'utilisateur cible
est l'ingénieur de validation BMS, qui passe déjà l'essentiel de sa
journée dans \textsc{MATLAB}.

L'application est implémentée comme une classe unique
\texttt{bms\_assistant} qui hérite de \texttt{matlab.apps.AppBase}.
Elle expose une seule méthode publique constructrice
(\texttt{bms\_assistant()}) et organise sa logique en méthodes
privées thématiques. Le thème graphique est volontairement clair et
sobre~: fond \texttt{[0,96~~0,97~~0,98]}, sidebar foncée
\texttt{[0,18~~0,22~~0,28]}, accents bleus \textsc{MATLAB}
\texttt{[0,00~~0,45~~0,74]} pour les actions primaires et orange
\texttt{[0,85~~0,33~~0,10]} pour les actions secondaires.

\subsection{Vue Dashboard}
\label{subsec:ch4_dashboard}

La vue \emph{Dashboard} est l'écran d'accueil. Elle affiche trois
informations~: le statut courant de l'outil (\emph{Ready}, \emph{Running},
\emph{OK}, \emph{Failed} ou \emph{Cancelled}), l'horodatage de la
dernière exécution, et le nombre de rapports archivés sur disque.
Elle propose en parallèle deux actions rapides~: lancer directement
une exécution avec les paramètres par défaut, et naviguer vers la
vue \emph{Reports} pour consulter les exécutions passées.

\begin{figure}[H]
  \centering
  \fbox{\parbox{0.85\textwidth}{\centering\vspace{2.2cm}
    \textit{[Capture vue Dashboard : sidebar à gauche avec 5 boutons
    nav (Dashboard surligné), zone principale avec 3 KPI cards (Status:
    Ready, Last run: 2026-05-11 04:05, Reports: 12 available) et 2
    boutons d'action rapide (Run validation, View reports)]}
    \vspace{2.2cm}}}
  \caption{Vue \emph{Dashboard} de l'application \texttt{bms\_assistant}.}
  \label{fig:ch4_ui_dashboard}
\end{figure}

\subsection{Vue Validate}
\label{subsec:ch4_validate}

La vue \emph{Validate} est le cœur opérationnel de l'application.
Elle expose les contrôles de configuration de la chaîne~: trois cases
à cocher pour activer ou désactiver MIL, SIL et PIL~; un bouton
permettant d'ouvrir une boîte de dialogue de sélection multiple des
scénarios MIL (\emph{checklist popup})~; deux listes déroulantes pour
sélectionner le scénario unique de SIL et de PIL~; et une case
\emph{Skip build} pour économiser le temps de génération de code C
lorsqu'aucune modification du SUT n'a eu lieu depuis le dernier
\emph{run}. Le bouton \emph{Run} déclenche
\texttt{orch.run} et active la barre de progression et les
\emph{phase chips}, qui changent de couleur (gris/bleu/vert) à
mesure que les phases s'exécutent. La console de \emph{logging} est
masquée par défaut et activée par un bouton \emph{Show log}, choix
ergonomique qui réduit l'encombrement visuel pendant les exécutions
\og{}sans incident\fg{}.

\begin{figure}[H]
  \centering
  \fbox{\parbox{0.9\textwidth}{\centering\vspace{2.5cm}
    \textit{[Capture vue Validate : sidebar surlignée sur Validate, en
    haut zone de configuration (3 checks MIL/SIL/PIL, scénarios MIL
    multi-sélection, dropdowns SIL et PIL, skip\_build), bouton Run
    central, barre de progression avec 6 chips colorées (mil/cov/sil/
    pil/facts/analyzer/doc), console toggle en bas, console déployée
    avec quelques lignes de log]}
    \vspace{2.5cm}}}
  \caption{Vue \emph{Validate} avec console déployée pendant une
  exécution MIL~+~SIL~+~PIL.}
  \label{fig:ch4_ui_validate}
\end{figure}

\subsection{Vue Reports}
\label{subsec:ch4_reports}

La vue \emph{Reports} liste les rapports HTML archivés sous
\texttt{doc\_generator/reports/} dans une table triable, et propose
quatre actions~: ouvrir un rapport dans le navigateur par défaut
(\emph{Open HTML}), exporter le \texttt{FULL\_REPORT.json}
correspondant (\emph{Save JSON}), exporter le rapport HTML lui-même
(\emph{Save HTML}), et supprimer un rapport (\emph{Delete}).
L'opération de suppression demande une confirmation et n'est jamais
silencieuse, conformément aux principes d'opérations destructrices
réversibles décrits au chapitre~\ref{chap:analyse_specification}.

\begin{figure}[H]
  \centering
  \fbox{\parbox{0.85\textwidth}{\centering\vspace{2.2cm}
    \textit{[Capture vue Reports : table avec colonnes timestamp,
    HTML file, JSON file, status, ligne sélectionnée surlignée, à
    droite barre d'actions verticale (Refresh, Open HTML, Save JSON,
    Save HTML, Delete avec icône poubelle rouge)]}
    \vspace{2.2cm}}}
  \caption{Vue \emph{Reports}~: consultation et gestion des rapports
  HTML archivés.}
  \label{fig:ch4_ui_reports}
\end{figure}

\subsection{Vue Plots}
\label{subsec:ch4_plots}

La vue \emph{Plots} affiche les figures PNG produites pendant la
dernière exécution de validation. Elle est articulée autour d'un
sélecteur de \emph{suite} (\emph{plant}, \emph{mil/bms},
\emph{mil/predictor}, \emph{sil}, \emph{pil}) et d'une liste des
fichiers PNG disponibles dans le sous-dossier correspondant de
\texttt{validator/reports/plots/}. La sélection d'un fichier dans la
liste affiche le PNG correspondant en taille réelle dans le panneau
principal, ce qui permet à l'ingénieur de revoir visuellement
n'importe quel test sans quitter l'application.

\begin{figure}[H]
  \centering
  \fbox{\parbox{0.85\textwidth}{\centering\vspace{2.2cm}
    \textit{[Capture vue Plots : sélecteur suite en haut, liste de
    PNG à gauche (PRD-DR-01.png surlignée), aperçu plein cadre du PNG
    sélectionné à droite (graphique signal MIL avec courbes simulation
    vs mesure)]}
    \vspace{2.2cm}}}
  \caption{Vue \emph{Plots}~: visualisation des figures PNG produites
  par les tests de validation.}
  \label{fig:ch4_ui_plots}
\end{figure}

\subsection{Vue Model}
\label{subsec:ch4_model}

La vue \emph{Model} expose les opérations sur les modèles
\textsc{Simulink} du SUT~: ouverture en édition pour le \emph{master},
le \emph{slave} et le \emph{predictor}, sauvegarde de la
configuration PIL au format \texttt{.c} pour archivage hors session
\textsc{MATLAB}, et réinitialisation du SUT à son état nominal par
restauration depuis \emph{snapshot}. Le mécanisme de \emph{snapshot}
est crucial~: à la première édition d'un modèle, une copie
\texttt{.bak} du fichier \texttt{.slx} est automatiquement créée
sous \texttt{model/.snapshots/}, ce qui garantit que l'utilisateur
peut toujours revenir à l'état initial après une session
exploratoire. La fonction \emph{Reset} parcourt l'ensemble des
\texttt{.bak} et restaure les fichiers correspondants après
confirmation explicite de l'utilisateur, conformément à la doctrine
opérationnelle de l'outil.

\begin{figure}[H]
  \centering
  \fbox{\parbox{0.85\textwidth}{\centering\vspace{2.2cm}
    \textit{[Capture vue Model : 5 boutons d'action verticaux (Edit
    master, Edit slave, Edit predictor, Save PIL .c, Reset to
    snapshot), zone de texte à droite listant l'état des modèles
    (OK / MISSING) et des snapshots]}
    \vspace{2.2cm}}}
  \caption{Vue \emph{Model}~: inspection, modification et restauration
  par \emph{snapshot} des modèles \textsc{Simulink}.}
  \label{fig:ch4_ui_model}
\end{figure}

\section{Validation unitaire du logiciel}
\label{sec:ch4_validation_unitaire}

\subsection{Stratégie de test et couverture visée}
\label{subsec:ch4_strategie_test}

La validation du logiciel suit la même règle d'or que celle du
modèle~: une exigence, un test. Les tests Python sont écrits avec
\texttt{pytest} et stockés sous \texttt{tests/python/}~; les tests
\textsc{MATLAB} sont écrits avec \texttt{matlab.unittest} et stockés
sous \texttt{tests/matlab/}. La cible de couverture est de 100~\%
sur les modules \texttt{doc\_generator} et \texttt{analyzer/python},
qui constituent le cœur métier du logiciel. Pour les modules à fort
couplage avec l'environnement (\texttt{orchestrator/cli.py} qui
appelle \texttt{matlab -batch}, \texttt{llm\_client.py} qui appelle
le réseau), la cible est ajustée à \emph{100~\% des chemins testables
sans dépendance externe}~: les appels HTTP et les sous-processus
\textsc{MATLAB} sont mockés (\texttt{monkeypatch} pour \texttt{pytest})
plutôt que réellement exercés.

\subsection{Tests unitaires — Module Validateur}
\label{subsec:ch4_tests_validateur}

Les tests \textsc{MATLAB} unitaires du Validateur sont stockés sous
\texttt{tests/matlab/}. Ils couvrent trois fichiers~:
\texttt{test\_orch\_list\_reports.m} qui valide l'énumération et le
tri des rapports archivés, \texttt{test\_orch\_scenarios.m} qui
vérifie la cohérence du catalogue de scénarios MIL/SIL/PIL exposé à
l'UI, et \texttt{test\_val\_new\_result\_check.m} qui couvre la
fabrique de résultats de tests (\texttt{val.new\_result} et
\texttt{val.check}). Les tests fonctionnels du Validateur sont, eux,
les tests MIL/SIL/PIL eux-mêmes décrits au chapitre~\ref{chap:partie_modele}~:
ils servent à la fois de validation du SUT et de validation du
Validateur dans la même exécution. Cette dualité est rendue possible
par le fait que le \emph{harness} \textsc{Simulink} et la sérialisation
JSON du verdict sont des composants logiciels indépendants du SUT.

\subsection{Tests unitaires — Module Analyseur}
\label{subsec:ch4_tests_analyseur}

Les tests \texttt{pytest} du module Analyseur se répartissent en
trois fichiers (\texttt{test\_facts.py}, \texttt{test\_rules.py},
\texttt{test\_engine.py}) et couvrent quatorze cas indépendants. Les
tests de \texttt{facts.py} valident l'assemblage du contexte de
travail à partir de fichiers JSON synthétiques, la résolution des
chemins pointés (présents et absents), et l'indexation des exigences
par identifiant. Les tests de \texttt{rules.py} valident le parsing
YAML, le rejet des règles malformées (identifiant manquant, premise
manquante), l'évaluation correcte d'une premise satisfaite et le
court-circuit silencieux d'une premise référençant un fait absent.
Les tests de \texttt{engine.py} couvrent les invariants
fonctionnels du moteur~: une règle ne déclenche qu'une fois par
exécution, l'expansion d'un \emph{forEach} produit un \emph{finding}
par binding, la logique de \emph{WIP downrank} rétrograde
correctement la sévérité, le \emph{locator} \texttt{hilite\_system}
est généré pour tout chemin contenant un \texttt{/}, et la sélection
de la pire sévérité dans \texttt{summarize} renvoie le rang attendu.

\subsection{Tests unitaires — Module Générateur de rapports}
\label{subsec:ch4_tests_doc_generator}

\begin{sloppypar}
Les tests du Générateur de rapports se répartissent en quatre
fichiers (\texttt{test\_consolidate.py}, \texttt{test\_prompt.py},
\texttt{test\_llm\_client.py}, \texttt{test\_cli\_doc.py}) et
couvrent vingt-deux cas indépendants. Les tests de
\texttt{consolidate.py} valident la forme du \emph{stub} de
validateur vide, le passage à l'identique des entrées dans le
\emph{full report}, le marquage des suites \emph{not\_run}, le calcul
correct des KPI à partir des suites présentes, l'élision de la
section \emph{analysis} dans le \emph{slim part 1}, et la
préservation de l'analyse dans le \emph{slim part 2}. Les tests de
\texttt{prompt.py} valident la présence des six identifiants de
section dans le \emph{prompt} de la Partie~1, des trois identifiants
dans le \emph{prompt} de la Partie~2, l'imposition du tag littéral
\texttt{<section>}, la rétro-compatibilité des alias, et la présence
du \emph{placeholder} \texttt{\{payload\}}. Les tests de
\texttt{llm\_client.py} valident l'envoi du \texttt{User-Agent}
explicite (cf. la note empirique sur le rejet Cloudflare), l'envoi
simultané des deux clés de \emph{tokens}, et le mécanisme de
\emph{retry} sur erreur 429. Les tests de \texttt{cli.py} (sept cas)
valident le \emph{trim} progressif sous budget, l'audit des
références JSON Pointer, le retrait des fences markdown, et la
distinction entre pointeurs et noms de fichiers.
\end{sloppypar}

\subsection{Tests fonctionnels — Orchestrateur}
\label{subsec:ch4_tests_orchestrateur}

Les tests Python de l'Orchestrateur (\texttt{test\_orchestrator.py})
sont au nombre de six et couvrent le contrat de l'\texttt{argparse}
exposé par \texttt{orchestrator/cli.py} ainsi que la composition
correcte du script \textsc{MATLAB} consommé par \texttt{matlab
-batch}. Les tests d'inclusion vérifient que les commandes
\texttt{mil.run}, \texttt{sil.run}, \texttt{pil.run},
\texttt{cov.run} et \texttt{ana.collect} apparaissent
\emph{sélectivement} dans le script en fonction des \emph{flags}
fournis, et que l'option \texttt{skip\_build} est correctement
propagée. Aucun sous-processus \texttt{matlab} réel n'est lancé~:
ces tests exercent uniquement la logique de composition.

\subsection{Tests de l'interface utilisateur}
\label{subsec:ch4_tests_ui}

L'interface utilisateur étant tributaire d'un display \textsc{MATLAB}
interactif, sa validation automatisée est limitée à l'interface
\emph{logique} entre l'UI et l'orchestrateur~: la vérification que la
liste de scénarios exposée par \texttt{orch.scenarios} est cohérente
avec celle attendue par \texttt{orch.run} (test
\texttt{test\_orch\_scenarios.m}), et que la fonction
\texttt{orch.list\_reports} renvoie une structure consommable par la
table \texttt{ReportsTable} (test \texttt{test\_orch\_list\_reports.m}).
La validation visuelle de l'UI fait l'objet d'une procédure
manuelle documentée dans
\texttt{tests/matlab/README.md} et exercée à chaque livraison~:
parcours des cinq vues, lancement d'une exécution complète,
vérification de la mise à jour des \emph{phase chips} et de la barre
de progression. Cette validation manuelle est consignée dans le
chapitre~\ref{chap:integration_validation} comme partie de la
recette d'intégration.

\subsection{Synthèse et taux de couverture atteint}
\label{subsec:ch4_synthese_tests}

La table~\ref{tab:ch4_synthese_tests} synthétise l'effort de test
et les taux de réussite de la dernière campagne consolidée.

\begin{table}[H]
  \centering
  \caption{Synthèse des tests unitaires du logiciel et de leurs
  résultats sur la dernière campagne (mai~2026).}
  \label{tab:ch4_synthese_tests}
  \small
  \begin{tabularx}{\textwidth}{|l|X|c|c|}
    \hline
    \textbf{Module}            & \textbf{Fichiers de test} & \textbf{Cas} & \textbf{Réussite} \\
    \hline
    Validateur (MATLAB)        & \texttt{test\_orch\_list\_reports.m},
                                 \texttt{test\_orch\_scenarios.m},
                                 \texttt{test\_val\_new\_result\_check.m} & 3 & 3 / 3 \\
    \hline
    Analyseur (Python)         & \texttt{test\_facts.py},
                                 \texttt{test\_rules.py},
                                 \texttt{test\_engine.py}                 & 14 & 14 / 14 \\
    \hline
    Générateur (Python)        & \texttt{test\_consolidate.py},
                                 \texttt{test\_prompt.py},
                                 \texttt{test\_llm\_client.py},
                                 \texttt{test\_cli\_doc.py}               & 22 & 22 / 22 \\
    \hline
    Orchestrateur (Python)     & \texttt{test\_orchestrator.py}           & 6 & 6 / 6 \\
    \hline
    \textbf{Total}             & 8 fichiers Python + 3 fichiers MATLAB    & \textbf{45} & \textbf{45 / 45} \\
    \hline
  \end{tabularx}
\end{table}

Les vingt-neuf exigences logicielles \texttt{REQ-SW-*} sont
couvertes à 100~\% par cette suite~: chaque exigence est associée à
au moins un test \emph{pytest} ou \emph{matlab.unittest} qui en
vérifie un invariant fonctionnel. La traçabilité exigence~$\to$~test
est maintenue dans \texttt{validator/requirements/requirements.json}
et croisée automatiquement par le script de \emph{coverage} de la
suite Python.

\section{Conclusion}
\label{sec:ch4_conclusion}

Ce chapitre a déroulé la branche Logicielle du V-model dual du
projet, en parallèle exact de la branche Modèle traitée au
chapitre~\ref{chap:partie_modele}. L'architecture modulaire en cinq
composants — Validateur, Analyseur, Générateur de rapports,
Orchestrateur et Interface utilisateur — a été conçue pour découpler
les responsabilités, isoler les choix technologiques externes
(notamment l'API Groq) et garantir la traçabilité de bout en bout
des artefacts produits.

Chaque module a été développé en suivant la même discipline~: une
seule source de vérité par information (paramètres dans
\texttt{init\_system.m}, règles dans
\texttt{analyzer/knowledge/}, scénarios dans
\texttt{orch.scenarios}), une sérialisation JSON systématique des
sorties pour permettre l'inspection et la comparaison, et une
gestion d'erreurs \emph{fail-fast} qui privilégie la lisibilité du
diagnostic sur la sophistication de la récupération. L'interface
graphique, développée en \textsc{MATLAB App~Designer}, expose
l'ensemble en cinq vues thématiques sous un style sobre et
professionnel cohérent avec l'environnement de travail quotidien de
l'ingénieur de validation BMS.

La validation unitaire du logiciel a porté sur quarante-cinq cas de
test (quarante-deux \texttt{pytest} + trois \texttt{matlab.unittest})
qui couvrent l'intégralité des vingt-neuf exigences
\texttt{REQ-SW-*}. La dernière campagne enregistre 45/45 succès, ce
qui clôt la phase de validation unitaire de la branche
Logicielle.

Le chapitre~\ref{chap:integration_validation} traite à présent la
phase \emph{commune} aux deux branches du V~: l'intégration des deux
sorties (système sous test du chapitre~3 et outil logiciel du
chapitre~4) en un système complet, et la phase d'acceptation qui
vérifie l'adéquation du tout aux besoins de l'utilisateur
ingénieur, par déroulé des scénarios documentés au
chapitre~\ref{chap:analyse_specification}.
